\documentclass[../Bachelor_LennartKoliwer.tex]{subfiles}
\graphicspath{{\subfix{../images/}}}
\begin{document}

\begin{lemma}\cite[43]{BaiSilverstein}\label{lem:AnzahlIsomklasse}
    Die Anzahl der Graphen in jeder Isomorphismenklasse kanonischer $\Delta(k,r,s)$ Graphen, mit festen $k$, $r$ und $s$ ist
    \begin{align*}
        \alpha = p(p-1)\cdots(p-r)n(n-1)\cdots(n-s+1)=p^{r+1}n^s\lr{1+\mathcal{O}\lr{\frac1n}}\,.
    \end{align*}
\end{lemma}

\begin{lemma}\cite[43]{BaiSilverstein}\label{lem:AnzahlEinezlKantenKateg3}
    Die Anzahl nicht doppelt belegter Kanten von $\Delta_3(k,r,s)$ Graphen ist kleiner gleich $k$.
\end{lemma}

\begin{lemma}\cite[43]{BaiSilverstein}\label{AnzahlKateg1Graphen}
    Die Anzahl der $\Delta_1(k,r,s)$ Graphen für festes k und r beträgt
    \begin{align*}
        \frac{1}{r+1}\binom{k}{r}\binom{k-1}{r}
    \end{align*}
\end{lemma}



\end{document}